\documentclass[a4paper,12pt]{article}

\usepackage[T1]{fontenc}
\usepackage{polyglossia}
\setmainlanguage{arabic}
\setotherlanguages{italian,english}

% Use a high-quality Arabic font available in most TeX distributions or Windows
\newfontfamily\arabicfont[Script=Arabic]{Amiri}
\newfontfamily\arabicfontsf[Script=Arabic]{Amiri}

\usepackage{graphicx}
\usepackage{float}
\usepackage{subcaption}
\usepackage{geometry}
\usepackage{fancyhdr}
\usepackage{xcolor}
\usepackage{mdframed}
\usepackage{enumitem}
\usepackage{hyperref}
\usepackage{booktabs}
\usepackage{array}
\usepackage{multicol}
\usepackage{bidi} % Essential for RTL support in XeLaTeX

% ── ألوان العلامة التجارية ─────────────────────────────────────────────────────────────
\definecolor{nutrigreen}{RGB}{46,125,50}
\definecolor{nutrilightgreen}{RGB}{232,245,233}
\definecolor{nutridark}{RGB}{28,28,28}
\definecolor{nutrigray}{RGB}{117,117,117}
\definecolor{nutriaccent}{RGB}{56,142,60}

% ── هوامش الصفحة ────────────────────────────────────────────────────────────
\setlength{\headheight}{30pt}
\geometry{left=2.5cm, right=2.5cm, top=2.8cm, bottom=2.5cm}

% ── رأس الصفحة ─────────────────────────────────────────────────────────────
\pagestyle{fancy}
\fancyhf{}
\fancyhead[L]{\small\color{nutrigray}القسم 4 --- دليل استخدام التطبيق}
\fancyhead[R]{\small\color{nutrigray}\thepage}
\renewcommand{\headrulewidth}{0.4pt}
\renewcommand{\headrule}{\hbox to\headwidth{\color{nutrigreen}\leaders\hrule height \headrulewidth\hfill}}

% ── عداد الخطوات ───────────────────────────────────────────────────────────
\newcounter{passo}[subsection]
\newcommand{\passaggio}[1]{%
    \stepcounter{passo}%
    \subsubsection*{\color{nutriaccent}الخطوة \thesubsection.\thepasso\ ---\ #1}%
}

% ── مربع معلومات أخضر ─────────────────────────────────────────────────────
\newmdenv[
  backgroundcolor=nutrilightgreen,
  linecolor=nutrigreen,
  linewidth=1.5pt,
  roundcorner=5pt,
  innertopmargin=8pt,
  innerbottommargin=8pt,
  innerleftmargin=12pt,
  innerrightmargin=12pt,
  skipabove=8pt,
  skipbelow=8pt
]{infobox}

% ── تسمية غامقة خضراء ─────────────────────────────────────────────────
\newcommand{\labelgreen}[1]{\medskip\noindent{\color{nutriaccent}\textbf{#1}}\par\smallskip}

% ── تلميح مائل ──────────────────────────────────────────────────────
\newcommand{\suggerimento}[1]{%
  \smallskip\noindent\textit{\color{nutrigray}#1}\par\smallskip%
}

% ── متغير لمجلد الصور (مفيد للترجمة) ───────────────────
\newcommand{\imgdir}{immagini_arb}

% ── صورة في المنتصف مع تعليق ─────────────────────────────────────────
\newcommand{\schermata}[3][0.38\textwidth]{%
  \begin{figure}[H]
    \centering
    \includegraphics[width=#1]{\imgdir/#2}
    \caption{\small\color{nutrigray}#3}
  \end{figure}%
}

% ── صورتان متجاورتان ──────────────────────────────────────────────────
\newcommand{\dueschermate}[4]{%
  \begin{figure}[H]
    \centering
    \begin{subfigure}[b]{0.42\textwidth}
      \centering
      \includegraphics[width=\textwidth]{\imgdir/#1}
      \caption{\small\color{nutrigray}#2}
    \end{subfigure}
    \hfill
    \begin{subfigure}[b]{0.42\textwidth}
      \centering
      \includegraphics[width=\textwidth]{\imgdir/#3}
      \caption{\small\color{nutrigray}#4}
    \end{subfigure}
  \end{figure}%
}

% ── ثلاث صور متجاورة ──────────────────────────────────────────────────
\newcommand{\treschermate}[6]{%
  \begin{figure}[H]
    \centering
    \begin{subfigure}[b]{0.30\textwidth}
      \centering
      \includegraphics[width=\textwidth]{\imgdir/#1}
      \caption{\small\color{nutrigray}#2}
    \end{subfigure}
    \hfill
    \begin{subfigure}[b]{0.30\textwidth}
      \centering
      \includegraphics[width=\textwidth]{\imgdir/#3}
      \caption{\small\color{nutrigray}#4}
    \end{subfigure}
    \hfill
    \begin{subfigure}[b]{0.30\textwidth}
      \centering
      \includegraphics[width=\textwidth]{\imgdir/#5}
      \caption{\small\color{nutrigray}#6}
    \end{subfigure}
  \end{figure}%
}

% ─────────────────────────────────────────────────────────────────────────────
\begin{document}

% ════════════════════════════════════════════════════════════════════════════
%  عنوان القسم
% ════════════════════════════════════════════════════════════════════════════
\begin{center}
  {\Huge\bfseries\color{nutrigreen} القسم 4}\\[6pt]
  {\Large\bfseries\color{nutridark} دليل استخدام التطبيق}\\[4pt]
  {\normalsize\itshape\color{nutrigray} دليل المستخدم --- تعليمات خطوة بخطوة}
\end{center}

\medskip
\noindent\rule{\textwidth}{1.5pt}{\color{nutrigreen}}
\medskip

\noindent يوضح هذا الفصل الميزات الرئيسية لتطبيق \textbf{NutriApp}
من خلال تعليمات مفصلة ولقطات شاشة مأخوذة مباشرة من التطبيق.
يتم تقسيم كل قسم إلى خطوات متتالية لتسهيل الفهم حتى
للمستخدمين الأقل خبرة.

\medskip
\begin{infobox}
  \textbf{\color{nutrigreen}فهرس الأقسام}
  \begin{enumerate}[leftmargin=1.4em, itemsep=2pt]
    \item تسجيل الدخول والتسجيل
    \item شريط التنقل
    \item إضافة الطعام
    \item تحديد الأهداف
    \item عرض التقدم
    \item إعدادات الحساب
  \end{enumerate}
\end{infobox}

\newpage

% ════════════════════════════════════════════════════════════════════════════
%  1. تسجيل الدخول والتسجيل
% ════════════════════════════════════════════════════════════════════════════
\section{تسجيل الدخول والتسجيل}

عند فتح التطبيق لأول مرة، يتم عرض شاشة الترحيب،
والتي يمكنك من خلالها الوصول إلى حساب موجود أو إنشاء حساب جديد.
يدعم NutriApp ثلاث طرق لتسجيل الدخول: البريد الإلكتروني وكلمة المرور،
والمصادقة عبر حساب Google، والمصادقة عبر Apple ID.
يمكنك أيضًا استكشاف التطبيق في وضع الضيف دون إنشاء حساب.

% ─── 1.1 ─────────────────────────────────────────────────────────────────────
\subsection{شاشة الترحيب وتسجيل الدخول}

\passaggio{افتح التطبيق}

\schermata[0.40\textwidth]{login.jpeg}{شاشة الترحيب لتطبيق \textbf{NutriApp}:
  في الأعلى يوجد الشعار الأخضر والعنوان الفرعي \emph{سجل الدخول لإدارة وجباتك،
  والمغذيات الكبرى والصغرى كلها في تطبيق واحد.} تعرض اللوحة البيضاء المركزية
  حقلي \textbf{البريد الإلكتروني} و \textbf{كلمة المرور} (مع أيقونة عين
  لإظهار/إخفاء الأحرف)، والزر الأخضر \textbf{LOG IN}، و
  رابط \emph{المتابعة بدون تسجيل الدخول}. تحت الفاصل \emph{أو} توجد
  أزرار \textbf{تسجيل الدخول باستخدام Google} و \textbf{تسجيل الدخول باستخدام Apple}، وفي الأسفل
  يوجد الرابط \emph{ليس لديك حساب؟ سجل الآن}.}

\labelgreen{الوصف}
عند فتح NutriApp لأول مرة، تظهر شاشة تسجيل الدخول باسم التطبيق
باللون الأخضر الداكن وعنوان فرعي يصف وظيفته بإيجاز.
تحتوي لوحة تسجيل الدخول المركزية على حقلين:

\begin{itemize}[leftmargin=1.4em, itemsep=2pt]
  \item \textbf{البريد الإلكتروني} --- عنوان البريد الإلكتروني المرتبط بالحساب،
        والمشار إليه بأيقونة المغلف على يسار الحقل
  \item \textbf{كلمة المرور} --- كلمة مرور الحساب؛ المس أيقونة العين المشطوبة
        على يمين الحقل للتبديل بين إظهار الأحرف
        كنص عادي أو إخفائها
\end{itemize}

بعد ملء كلا الحقلين، اضغط على الزر الأخضر \textbf{تسجيل الدخول}
لإكمال المصادقة والدخول إلى التطبيق.

\labelgreen{المتابعة بدون تسجيل الدخول}
من خلال النقر على رابط \emph{المتابعة بدون تسجيل الدخول}، يمكنك استكشاف
التطبيق في وضع الضيف، دون إنشاء أو إدخال بيانات الاعتماد.
في هذا الوضع، لا تتم مزامنة البيانات مع السحابة ولا يمكن
الوصول إليها من أجهزة أخرى. يوصى بإنشاء حساب حتى لا
تفقد البيانات التي تم إدخالها.

\suggerimento{إذا نسيت كلمة المرور الخاصة بك، يمكنك استعادتها عبر
إجراء الاسترداد المتاح على الموقع الرسمي لـ NutriApp: يرسل
النظام رابط إعادة تعيين إلى عنوان البريد الإلكتروني المرتبط بالحساب.}

% ─── 1.2 ─────────────────────────────────────────────────────────────────────
\subsection{تسجيل الدخول باستخدام Google}

\passaggio{حدد تسجيل الدخول السريع عبر Google}

\dueschermate{login.jpeg}{شاشة تسجيل الدخول: في أسفل اللوحة البيضاء يوجد
  زر \textbf{تسجيل الدخول باستخدام Google} مع شعار G على اليسار. بالنقر عليه
  يفتح محدد حسابات Google الموجودة على الجهاز.}{login_google.jpeg}{%
  نافذة النظام \textbf{اختيار حساب}: تعرض جميع حسابات Google
  المحفوظة على الجهاز مع صورة الملف الشخصي وعنوان البريد الإلكتروني (مخفية
  للخصوصية). في الأسفل يوجد خيار \emph{إضافة حساب آخر} لإضافة
  حساب غير موجود بعد، والملاحظة التي تفيد بأن Google ستشارك اسمك،
  وبريدك الإلكتروني، وصورة ملفك الشخصي مع NutriApp.}

\labelgreen{الوصف}
انقر على \textbf{تسجيل الدخول باستخدام Google} في شاشة تسجيل الدخول لبدء
المصادقة السريعة. يفتح نظام التشغيل تلقائيًا نافذة تحديد حساب Google،
والتي تعرض جميع حسابات Gmail التي تم تكوينها بالفعل
على الجهاز. يُظهر كل إدخال في القائمة صورة الملف الشخصي
وعنوان البريد الإلكتروني للحساب.

\labelgreen{كيفية المتابعة}
\begin{itemize}[leftmargin=1.4em, itemsep=2pt]
  \item \textbf{انقر على حسابك} في القائمة لتحديده و
        تفويض NutriApp للوصول إلى اسمك وبريدك الإلكتروني وصورة ملفك الشخصي.
        يكتمل تسجيل الدخول تلقائيًا دون خطوات إضافية.
  \item \textbf{انقر على \emph{إضافة حساب آخر}} إذا كان الحساب المطلوب
        غير موجود في القائمة: سيتم فتح تدفق تسجيل الدخول القياسي لـ Google
        لإضافة حساب جديد إلى الجهاز.
\end{itemize}

\labelgreen{الخطوة التالية}
بعد تحديد الحساب، يصادق التطبيق المستخدم وينتقل
مباشرة إلى الشاشة الرئيسية. إذا كان هذا هو أول تسجيل دخول باستخدام حساب Google هذا،
فيبدأ تدفق الإعداد الأولي للملف الشخصي.

\suggerimento{يعد تسجيل الدخول باستخدام Google أسرع طريقة: فهو لا يتطلب
منك تذكر كلمة مرور منفصلة ويتم تحديثه تلقائيًا إذا قمت بتغيير
كلمة مرور Google الخاصة بك. يوصى به لأولئك الذين يستخدمون Google بالفعل في خدمات أخرى.}

% ─── 1.3 ─────────────────────────────────────────────────────────────────────
\subsection{إنشاء حساب جديد}

\passaggio{افتح نموذج التسجيل --- الخطوة 1}

\schermata[0.40\textwidth]{registra_step1.jpeg}{شاشة \textbf{التسجيل}
  --- الخطوة 1 من 3: يُظهر مؤشر التقدم في الأعلى الشرطة الأولى
  باللون الأخضر (نشطة) واثنتين باللون الرمادي (للاكتساب). في المنتصف توجد الصورة الرمزية المؤقتة مع
  أيقونة الكاميرا الخضراء لإضافة صورة اختيارية. حقول النموذج هي:
  \textbf{الاسم}، \textbf{اللقب}، \textbf{البريد الإلكتروني}، \textbf{كلمة المرور}، و
  \textbf{تأكيد كلمة المرور}. في الأسفل يوجد الزر الأخضر \textbf{SIGN UP}.}

\labelgreen{الوصف}
انقر على الرابط الأخضر \emph{ليس لديك حساب؟ سجل الآن} في شاشة تسجيل الدخول
لفتح معالج التسجيل المكون من ثلاث خطوات، والمعروض
بواسطة مؤشر التقدم في الأعلى (شرطة خضراء = الخطوة الحالية،
شرطات رمادية = الخطوات القادمة).

\textbf{الخطوة 1 --- البيانات الشخصية وبيانات الاعتماد.}
املأ الحقول التالية:

\begin{itemize}[leftmargin=1.4em, itemsep=2pt]
  \item \textbf{الاسم} و \textbf{اللقب} --- الاسم الكامل الذي سيتم
        عرضه في التطبيق وفي الترحيب على الشاشة الرئيسية
  \item \textbf{البريد الإلكتروني} --- عنوان البريد الإلكتروني الذي سيكون بمثابة
        المعرف الفريد للحساب؛ يجب أن يكون صالحًا لأنه
        سيتم استخدامه للتحقق عبر رمز OTP في الخطوة التالية
  \item \textbf{كلمة المرور} --- اختر كلمة مرور قوية تتكون من 8
        أحرف على الأقل؛ يوصى بتضمين أحرف كبيرة وصغيرة
        وأرقام وأحرف خاصة
  \item \textbf{تأكيد كلمة المرور} --- أعد إدخال نفس كلمة المرور
        لمنع الأخطاء المطبعية؛ يجب أن يتطابق الحقلان
        تمامًا للمتابعة
\end{itemize}

\textbf{صورة الملف الشخصي (اختياري).}
انقر على أيقونة الكاميرا الخضراء المتراكبة على الصورة الرمزية الرمادية لالتقاط
صورة أو اختيار واحدة من المعرض. الصورة اختيارية ويمكن
إضافتها أو تغييرها في أي وقت لاحق من
قسم \emph{ملفي الشخصي} في الإعدادات.

\labelgreen{الخطوة التالية}
تأكد من ملء جميع الحقول الإلزامية بشكل صحيح،
ثم اضغط على الزر الأخضر \textbf{SIGN UP} للانتقال إلى الخطوة
التالية.

\passaggio{تحقق من عنوان بريدك الإلكتروني --- الخطوة 2}

\schermata[0.40\textwidth]{registra_step2.jpeg}{شاشة \textbf{التسجيل}
  --- الخطوة 2 من 3: يُظهر المؤشر أول شرطتين باللون الأخضر والأخيرة
  باللون الرمادي. في المنتصف توجد أيقونة مغلف بعلامة اختيار خضراء وعنوان
  \textbf{تحقق من بريدك الإلكتروني}. يشير النص إلى أنه
  تم إرسال رمز OTP إلى عنوان البريد الإلكتروني الذي تم إدخاله (مخفي جزئيًا).
  ينتظر حقل \textbf{رمز OTP} الإدخال. يسمح لك الرابط
  \emph{إرسال مرة أخرى} بطلب رمز جديد.
  في الأسفل يوجد الزر الأخضر \textbf{CHECK EMAIL}.}

\labelgreen{الوصف}
بعد الضغط على \textbf{متابعة} في الخطوة الأولى، يرسل التطبيق تلقائيًا
رمز OTP (كلمة مرور لمرة واحدة) إلى عنوان البريد الإلكتروني الذي تم إدخاله.
تعرض الشاشة أيقونة مغلف بعلامة اختيار خضراء و
تشير إلى العنوان الذي تم إرسال الرمز إليه.

\labelgreen{كيفية إكمال التحقق}
\begin{itemize}[leftmargin=1.4em, itemsep=2pt]
  \item افتح صندوق الوارد لعنوان البريد الإلكتروني المستخدم أثناء
        التسجيل وابحث عن البريد الإلكتروني للتحقق من NutriApp
  \item لاحظ أو انسخ رمز OTP المكون من 6 أرقام الموجود في البريد الإلكتروني
  \item عد إلى التطبيق وأدخل الرمز في حقل \textbf{رمز OTP}
  \item اضغط على \textbf{CHECK EMAIL} لتأكيد عنوان البريد الإلكتروني
        والانتقال إلى الخطوة الأخيرة
\end{itemize}

\labelgreen{لم يتم استلام الرمز}
إذا لم يصل الرمز في غضون بضع دقائق، فتحقق من مجلد
\emph{البريد العشوائي} أو \emph{البريد غير الهام} لدى مزودك. بدلاً من ذلك،
انقر على الرابط \emph{إرسال مرة أخرى}
لطلب رمز OTP جديد. يصبح الرابط نشطًا بعد فترة
انتظار قصيرة لمنع الطلبات المتعددة السريعة.

\suggerimento{رمز OTP صالح لفترة زمنية محدودة (عادة
10--15 دقيقة). إذا انتظرت طويلاً قبل إدخاله،
ستحتاج إلى طلب رمز جديد عبر رابط \emph{إرسال مرة أخرى}.}

\passaggio{إكمال التسجيل --- الخطوة 3}

\schermata[0.40\textwidth]{registra_step3.jpeg}{شاشة \textbf{التسجيل}
  --- الخطوة 3 من 3: جميع الشرطات الثلاث الموجودة على المؤشر خضراء، مما يشير إلى
  اكتمال العملية. في المنتصف توجد دائرة خضراء بها علامة اختيار بيضاء و
  عنوان \textbf{عمل رائع!} مع عنوان فرعي \emph{أنت جاهز لـ
  بدء رحلتك مع NutriApp.} في الأسفل يوجد الزر الأخضر
  \textbf{GO TO HOMEPAGE}.}

\labelgreen{الوصف}
بعد التحقق الناجح من رمز OTP، يعرض التطبيق شاشة التأكيد.
تتحول جميع شرطات المؤشر الثلاث إلى اللون الأخضر. تؤكد الرسالة
\textbf{عمل رائع!} أنه تم إنشاء الحساب و
التحقق منه بشكل صحيح. اضغط على \textbf{GO TO HOMEPAGE} للدخول
إلى التطبيق.

\begin{infobox}
  \textbf{\color{nutrigreen}ملخص طرق تسجيل الدخول}

  \smallskip
  \begin{tabular}{@{}lll@{}}
    \toprule
    \textbf{الطريقة}       & \textbf{المتطلبات}             & \textbf{مثالي لـ} \\
    \midrule
    البريد الإلكتروني + كلمة المرور      & حساب NutriApp                 & أولئك الذين يفضلون بيانات اعتماد مخصصة \\
    تسجيل الدخول باستخدام Google    & حساب Gmail على الجهاز          & وصول سريع بدون كلمة مرور \\
    تسجيل الدخول باستخدام Apple     & Apple ID (iOS فقط)              & مستخدمو iPhone/iPad \\
    بدون تسجيل الدخول    & لا شيء                             & استكشاف سريع للتطبيق \\
    \bottomrule
  \end{tabular}
\end{infobox}

\newpage

% ════════════════════════════════════════════════════════════════════════════
%  2. شريط التنقل
% ════════════════════════════════════════════════════════════════════════════
\section{شريط التنقل}

شريط التنقل هو العنصر الثابت للواجهة: يظهر في
أسفل كل شاشة ويسمح لك بالانتقال فورًا بين
المجالات الرئيسية الخمسة للتطبيق بنقرة واحدة.

% ─── 2.1 ─────────────────────────────────────────────────────────────────────
\subsection{عرض شريط التنقل}

\passaggio{افتح التطبيق}

\schermata[0.38\textwidth]{homepage.jpg}{الشاشة الرئيسية \textbf{NutriApp}:
  ملخص يومي للسعرات الحرارية مع رسم بياني دائري، وأشرطة المغذيات الكبرى
  (الكربوهيدرات، البروتين، الدهون) والفترات الزمنية الأربع للوجبات. شريط التنقل
  في الأسفل بأيقونات NutriApp، والتقويم، والإضافة، والوصفات،
  والرسوم البيانية.}

\labelgreen{الوصف}
عند فتح التطبيق، يتم عرض الشاشة الرئيسية مع
تحية مخصصة في الأعلى، والتاريخ الحالي القابل للتنقل باستخدام السهمين \texttt{<} و \texttt{>} ،
ورسم بياني دائري أخضر يشير إلى السعرات الحرارية المتبقية لهذا اليوم.
إلى يمين الرسم البياني توجد أشرطة ملونة أفقية للكربوهيدرات (البرتقالي)،
والبروتين (الأخضر)، والدهون (الوردي/الأحمر)، يعرض كل منها القيمة المستهلكة مقارنة بالهدف.
الفترات الزمنية الأربع --- الإفطار، والغداء، والعشاء، والوجبات الخفيفة --- هي
يمكن الوصول إليها عبر زر \textbf{+} بجوار كل منها.

\labelgreen{عناصر الشريط}
\begin{itemize}[leftmargin=1.4em, itemsep=2pt]
  \item \textbf{NutriApp} (أيقونة المنزل) --- الشاشة الرئيسية مع ملخص
        السعرات الحرارية والمغذيات الكبرى وأقسام الوجبات اليومية
  \item \textbf{التقويم} (أيقونة التقويم) --- عرض شهري برمز لون أخضر/أحمر
        للأيام مع أو بدون وجبات مسجلة
  \item \textbf{إضافة} (الزر الأخضر المركزي \textbf{+}) --- وصول سريع
        إلى شاشة إدخال الطعام
  \item \textbf{الوصفات} (أيقونة كتاب) --- قائمة الوصفات المحفوظة
  \item \textbf{الرسوم البيانية} (أيقونة رسم بياني خطي) --- إحصائيات التغذية
        مع طرق عرض يومية وأسبوعية وشهرية وسنوية
\end{itemize}

\suggerimento{يتم تمييز الأيقونة النشطة باللون الأخضر في الشريط السفلي.
يتوفر الزر الأخضر \textbf{+} في المنتصف دائمًا لإضافة
عنصر غذائي بسرعة: إنه أسرع طريق للتسجيل اليومي.}

\newpage

% ════════════════════════════════════════════════════════════════════════════
%  3. إضافة الطعام
% ════════════════════════════════════════════════════════════════════════════
\section{إضافة الطعام}

يوفر التطبيق طرقًا متعددة لتسجيل الوجبات: البحث في قاعدة البيانات العالمية،
واختيار من الأطعمة المحفوظة، وإضافة وصفة تم إنشاؤها بالفعل، و
مسح الرمز الشريطي للمنتج.

% ─── 3.1 ─────────────────────────────────────────────────────────────────────
\subsection{افتح شاشة الإدخال}

\passaggio{انتقل إلى الشاشة الرئيسية}

\schermata[0.38\textwidth]{homepage.jpg}{الشاشة الرئيسية: رسم بياني دائري مع السعرات
  المتبقية، وأشرطة المغذيات الكبرى، والفترات الزمنية الأربع التي يحتوي كل منها على
  زر \textbf{+}.}

\labelgreen{الخطوة التالية}
اضغط على \textbf{+} بجوار الفترة الزمنية المطلوبة (الإفطار، الغداء،
العشاء، أو وجبة خفيفة) أو اضغط على الزر الأخضر الكبير \textbf{+} في وسط
شريط التنقل للوصول إلى شاشة \emph{إضافة}.

% ─── 3.2 ─────────────────────────────────────────────────────────────────────
\subsection{اختر طريقة الإدخال}

\passaggio{حدد مصدر الطعام}

\dueschermate{pasto_dropdown.jpg}{شاشة \textbf{يدوي}: قائمة منسدلة
  \emph{اختيار وجبة} وحقول لاسم الطعام، الوزن (جم)، والمغذيات الكبرى
  (السعرات الحرارية، الكربوهيدرات، البروتين، الدهون، الألياف، السكريات).}{pasto_recenti.jpg}{%
  شاشة \textbf{إضافة} مع تنشيط علامة تبويب \emph{إضافة طعام}: شريط بحث
  مع أيقونة باركود. يوجد أدناه رابط \emph{لا يمكنك العثور على الطعام؟
  سجله بنفسك!}}

\labelgreen{الوصف}
تحتوي شاشة \emph{إضافة} على ثلاث علامات تبويب في الأعلى:

\begin{itemize}[leftmargin=1.4em, itemsep=2pt]
  \item \textbf{إضافة طعام} --- البحث في قاعدة البيانات العالمية أو مسح
        الرمز الشريطي
  \item \textbf{الوصفات المحفوظة} --- الوصفات التي أنشأها المستخدم مع إجمالي
        السعرات الحرارية لكل منها
  \item \textbf{الأطعمة المحفوظة} --- الأطعمة التي تم تسجيلها يدويًا سابقًا
\end{itemize}

\passaggio{التنقل بين علامات تبويب الوصفات المحفوظة والأطعمة المحفوظة}

\dueschermate{pasto_ricette.jpg}{علامة تبويب \textbf{الوصفات المحفوظة}: تُظهر كل وصفة
  اسمها وإجمالي السعرات الحرارية. في الأسفل يوجد رابط
  \emph{إدراج وصفة جديدة}.}{pasto_preferiti.jpg}{علامة تبويب \textbf{الأطعمة المحفوظة}:
  شريط البحث؛ إذا لم تكن هناك أطعمة محفوظة، تظهر الرسالة
  \emph{لا يوجد طعام محفوظ}.}

\suggerimento{يتم حفظ الأطعمة التي تم إدخالها يدويًا تلقائيًا
في قسم \emph{الأطعمة المحفوظة} وستكون متاحة للإدخالات المستقبلية
دون الحاجة إلى إعادة إدخالها من الصفر.}

% ─── 3.3 ─────────────────────────────────────────────────────────────────────
\subsection{إدخال الطعام يدويًا}

\passaggio{املأ نموذج الإدخال}

\dueschermate{registra_alimento.jpg}{نموذج \textbf{المكون}: حقول
  \emph{اسم الطعام} و \emph{الوزن (جم)}، متبوعًا بـ
  قسم \textbf{المغذيات الكبرى} بالكربوهيدرات والبروتين والدهون والألياف،
  والسكريات، والماء.}{nutrienti_micro1.jpg}{قسم \textbf{المكون} بأقسام قابلة للتوسيع
  \emph{تفاصيل الدهون}، \emph{الفيتامينات}، و
  \emph{المعادن} لتتبع غذائي كامل.}

\labelgreen{الوصف}
انقر على \emph{سجله بنفسك!} لفتح نموذج الإدخال اليدوي.
املأ على الأقل الحقول الإلزامية في قسم \textbf{المغذيات الكبرى} عن طريق
قراءة القيم من عبوة المنتج (ملصق \emph{القيم الغذائية
لكل 100 جم}). للتتبع الكامل، يمكنك توسيع
أقسام \textbf{تفاصيل الدهون} و \textbf{الفيتامينات} و \textbf{المعادن}.


% ─── 3.4 ─────────────────────────────────────────────────────────────────────
\subsection{الإدخال عبر مسح الرمز الشريطي (الباركود)}

\passaggio{بدء كاميرا المسح الضوئي}

\schermata[0.38\textwidth]{scansione_barcode.jpg}{شاشة المسح
  \textbf{الباركود}: تركز الكاميرا على الجزء الخلفي من العبوة.
  يحدد الإطار الأخضر منطقة التعرف. في الجزء العلوي توجد عناصر التحكم
  \textbf{تشغيل الفلاش} و \textbf{إلغاء}.}

\labelgreen{الوصف}
انقر على أيقونة الباركود الموجودة على يمين شريط البحث في
علامة التبويب \emph{إضافة طعام}. ضع الباركود داخل
المربع الأخضر: يحدث التعرف تلقائيًا بمجرد أن
يصبح الرمز في بؤرة التركيز، دون الضغط على أي أزرار.

\suggerimento{في البيئات ضعيفة الإضاءة، قم بتنشيط الفلاش بالنقر على
\textbf{تشغيل الفلاش} في أعلى اليمين. احمل الهاتف على بعد حوالي 15--20 سم
من العبوة للحصول على أفضل النتائج.}

% ─── 3.5 ─────────────────────────────────────────────────────────────────────
\subsection{إنشاء وصفة جديدة}

\passaggio{افتح النموذج وأضف المكونات}

\treschermate{nuova_ricetta.jpg}{نموذج \textbf{إنشاء وصفة}: حقول
  \emph{اسم الوصفة}، \emph{الكميات}، و \emph{الإجراء / الملاحظات}.
  يُظهر قسم المكونات \emph{تمت إضافة 0}. في الأسفل توجد أزرار
  \textbf{+ مكونات} و \textbf{حفظ}.}{nuova_ricetta_recenti.jpg}{%
  شاشة \textbf{إضافة مكون} مع علامة تبويب \emph{إضافة طعام}:
  شريط بحث مع باركود. أدناه الرابط
  \emph{لا يمكنك العثور على الطعام؟ سجله بنفسك!}}{nuova_ricetta_preferiti.jpg}{%
  علامة التبويب \emph{الأطعمة المحفوظة}: قائمة بالأطعمة المسجلة بالفعل أو
  رسالة \emph{لا يوجد طعام محفوظ} إذا كانت فارغة.}

\labelgreen{الوصف}
من قسم \emph{الوصفات}، انقر على الزر الأخضر \textbf{+} لفتح
نموذج \emph{إنشاء وصفة}. املأ الاسم، وعدد الحصص، والملاحظات الاختيارية،
ثم اضغط على \textbf{+ مكونات} لإضافة كل
مكون. يظهر كل مكون مضاف في القائمة مع أيقونات
التعديل والحذف. عند الانتهاء، اضغط على \textbf{حفظ}.

% ─── 3.6 ─────────────────────────────────────────────────────────────────────
\subsection{عرض تفاصيل الوجبة}

\passaggio{افتح شاشة تفاصيل الفترة الزمنية}

\schermata[0.38\textwidth]{dettaglia_pasto.jpg}{شاشة \textbf{الإفطار}: لافتة برتقالية
  مع أيقونة وعداد طعام، والتنقل بين التواريخ، وزر \textbf{+}
  لإضافة الأطعمة. يعرض \emph{الملخص الغذائي} أشرطة تقدم
  للسعرات الحرارية والكربوهيدرات والبروتين والدهون. في الأسفل توجد قائمة
  \emph{الأطعمة المستهلكة} مع أيقونات التعديل والحذف.}

\labelgreen{الوصف}
انقر فوق فترة زمنية على الشاشة الرئيسية لفتح الملخص التفصيلي
مع اللافتة الملونة، والملخص الغذائي بأشرطة التقدم،
وقائمة الأطعمة المستهلكة مع أيقونات التعديل (قلم رصاص أزرق)
والحذف (سلة مهملات حمراء).

\newpage

% ════════════════════════════════════════════════════════════════════════════
%  4. تحديد الأهداف
% ════════════════════════════════════════════════════════════════════════════
\section{تحديد الأهداف}

تحدد الأهداف القيم المرجعية للسعرات الحرارية والمغذيات الكبرى،
والوزن الذي يستخدمه التطبيق لحساب التقدم. تحديد أهداف واقعية
ضروري للحصول على بيانات تتبع ذات مغزى.

% ─── 4.1 ─────────────────────────────────────────────────────────────────────
\subsection{الوصول إلى الأهداف وتكوينها}

\passaggio{افتح الإعدادات وحدد أهدافي}

\dueschermate{impostazioni.jpg}{شاشة \textbf{الإعدادات}: أقسام للحساب
  (ملفي الشخصي، أهدافي، الإشعارات، اللغة و
  البلد)، الدعم (مساعدة، سياسة الخصوصية، معلومات عنا)، والإجراءات (الإبلاغ عن
  مشكلة، تسجيل الخروج).}{obiettivi.jpg}{%
  شاشة \textbf{أهدافي}: قسم \emph{الوزن} مع \emph{الوزن
  الحالي} (60.0 كجم) و \emph{الوزن المستهدف} (70.0 كجم). قسم
  \emph{المغذيات الكبرى} مع \emph{هدف السعرات الحرارية} (2000.0 كيلو كالوري)،
  \emph{الكربوهيدرات} (130.0 جم)، \emph{البروتين} (45.0 جم)، و
  \emph{الدهون} (36.0 جم). في الأسفل الأقسام القابلة للتوسيع
  \emph{تفاصيل الماكرو}، \emph{الفيتامينات}، و \emph{المعادن}.}

\labelgreen{الوصف}
انقر على أيقونة الترس في أعلى اليمين لفتح الإعدادات،
ثم حدد \textbf{أهدافي}. تجمع الشاشة المعلمات
في قسمين:

\begin{itemize}[leftmargin=1.4em, itemsep=2pt]
  \item \textbf{الوزن} --- \emph{الوزن الحالي} هو نقطة البداية لـ
        حساب الاحتياجات؛ \emph{الوزن المستهدف} هو الهدف
        المراد الوصول إليه (يمكن أن يكون أعلى أو أقل من الوزن الحالي)
  \item \textbf{المغذيات الكبرى} --- \emph{هدف السعرات الحرارية} يحدد الحد
        اليومي للسعرات؛ تحدد حقول الكربوهيدرات والبروتين والدهون
        الحدود بالجرام
\end{itemize}

\labelgreen{الخطوة التالية}
أدخل القيم المطلوبة واضغط على \textbf{حفظ}. يتم تطبيق التغييرات
على الفور في جميع أنحاء التطبيق. تظهر لافتة التأكيد الخضراء
في أسفل الشاشة.

\begin{infobox}
  \textbf{\color{nutrigreen}كيفية اختيار القيم الصحيحة:}
  نقطة البداية الشائعة هي الكربوهيدرات 45--65\%، والبروتين 10--35\%،
  الدهون 20--35\% من إجمالي السعرات الحرارية اليومية. استشر
  أخصائي التغذية للحصول على قيم مخصصة.
\end{infobox}

\newpage

% ════════════════════════════════════════════════════════════════════════════
%  5. عرض التقدم
% ════════════════════════════════════════════════════════════════════════════
\section{عرض التقدم}

يوفر التطبيق أربعة طرق عرض زمنية لتحليل السعرات الحرارية و
اتجاهات التغذية: يومية، أسبوعية، شهرية، وسنوية،
يمكن الوصول إليها من قسم \textbf{الرسوم البيانية} في شريط التنقل.

% ─── 5.1 ─────────────────────────────────────────────────────────────────────
\subsection{الوصول إلى بيانات التقدم}

\passaggio{افتح قسم الرسوم البيانية}

\schermata[0.38\textwidth]{storico_giornaliero.jpg}{شاشة \textbf{الرسوم البيانية}
  مع تنشيط علامة تبويب \emph{يومي}: مربع أخضر \emph{إحصائيات
  غذائية} يتضمن العناصر المسجلة والأيام المتتبعة وتاريخ آخر
  إدخال. يوجد أدناه محدد المرشح، وعلامات تبويب الوقت الأربعة، والمربع الأحمر
  مع إجمالي السعرات الحرارية.}

\labelgreen{الوصف}
انقر على أيقونة \textbf{الرسوم البيانية} في شريط التنقل. في الأعلى يظهر
المربع الأخضر مع الملخص العام، ثم لوحة \textbf{الفلاتر}
مع العناصر الغذائية المحددة (الافتراضي: السعرات الحرارية) والنطاق الزمني.
تتيح لك علامات التبويب الأربعة تغيير دقة العرض.

% ─── 5.2 ─────────────────────────────────────────────────────────────────────
\subsection{تحليل الرسوم البيانية}

\passaggio{مخطط شريطي --- السعرات الحرارية بمرور الوقت}

\schermata[0.40\textwidth]{grafici_barre.jpeg}{شاشة \textbf{الرسوم البيانية} مع
  تنشيط علامة تبويب \emph{يومي}. في الجزء العلوي، لوحة \textbf{تصدير تقرير PDF}
  مع زري \textbf{طباعة PDF} (أخضر) و \textbf{مشاركة PDF}. يُظهر مخطط
  \textbf{السعرات الحرارية بمرور الوقت} ثلاثة أشرطة عمودية حمراء لـ
  28 و 29 و 30 أبريل: 779 كيلو كالوري، 143 كيلو كالوري، و 377 كيلو كالوري. المقياس
  على المحور الرأسي يمتد من 0.0 إلى 779 كيلو كالوري. أدناه يبدأ قسم
  \emph{توزيع المغذيات الكبرى}.}

\labelgreen{الوصف}
يُظهر المخطط الشريطي لـ \textbf{السعرات الحرارية بمرور الوقت} السعرات
اليومية للفترة المحددة. يمثل كل شريط عمودي أحمر
يومًا واحدًا، مع الإشارة إلى القيمة بوحدة كيلو كالوري أعلاه. يتناسب الارتفاع
مع إجمالي السعرات الحرارية: الأيام التي تحتوي على عدد قليل من الأطعمة المسجلة تنتج أشرطة
أقصر (مثل 29 أبريل مع 143 كيلو كالوري)، والأيام التي بها وجبات أكثر تنتج
أشرطة أطول (مثل 28 أبريل مع 779 كيلو كالوري). يتكيف المحور الرأسي تلقائيًا
إلى أقصى قيمة للفترة.

\labelgreen{تصدير تقرير PDF}
فوق المخطط، توفر لوحة \textbf{تصدير تقرير PDF} زرين:
\begin{itemize}[leftmargin=1.4em, itemsep=2pt]
  \item \textbf{طباعة PDF} (أخضر) --- يولد المستند ويفتح نافذة الطباعة الخاصة بنظام
        التشغيل
  \item \textbf{مشاركة PDF} (أبيض بحدود) --- يولد المستند و
        يفتح لوحة المشاركة لإرسالها عبر البريد الإلكتروني أو WhatsApp أو
        أي تطبيق آخر مثبت
\end{itemize}

\passaggio{المخطط الدائري --- توزيع المغذيات الكبرى}

\schermata[0.40\textwidth]{grafici_torta.jpeg}{شاشة \textbf{الرسوم البيانية} ---
  قسم \textbf{توزيع المغذيات الكبرى}: مخطط حلقي يوضح
  النسبة المئوية للمغذيات الكبرى. تشير وسيلة الإيضاح إلى الكربوهيدرات
  44.7\%، والبروتين 12.0\%، والدهون 43.3\%. الأشرطة الأفقية أدناه
  تُظهر القيم المطلقة: الكربوهيدرات 155.6 جم (شريط برتقالي)،
  البروتين 41.9 جم (شريط أزرق)، والدهون 67.1 جم (شريط أرجواني).}

\labelgreen{الوصف}
يكشف التمرير لأسفل في شاشة \emph{الرسوم البيانية} عن قسم
\textbf{توزيع المغذيات الكبرى}، والذي يعرض تمثيلين متكاملين:

\begin{itemize}[leftmargin=1.4em, itemsep=2pt]
  \item \textbf{مخطط حلقي} --- يوضح بصريًا النسبة بين
        المغذيات الكبرى الثلاثة: يمثل الجزء الذهبي/البرتقالي
        الكربوهيدرات، والأزرق الداكن للبروتين، والأرجواني للدهون.
        تتيح لك النسب المئوية الموجودة في وسيلة الإيضاح على اليمين قراءة القيم
        الدقيقة في لمح البصر
  \item \textbf{أشرطة أفقية} --- تُظهر كلاً من النسبة المئوية و
        القيمة المطلقة بالجرام لكل عنصر من المغذيات الكبرى، مما يسهل
        مقارنتها بالأهداف المحددة في قسم \emph{أهدافي}
\end{itemize}

\labelgreen{كيفية تفسير البيانات}
في المثال الموضح، تشكل الدهون 43.3\% من إجمالي الطاقة
المستهلكة (67.1 جم)، وهي قيمة أعلى من الحد الموصى به
من 20--35\%. بلغت نسبة الكربوهيدرات 44.7\% (155.6 جم) والبروتين
عند 12.0\% (41.9 جم). تتيح لك هذه المعلومات تحديد الاختلالات
في توزيع المغذيات الكبرى وتعديل الخيارات الغذائية
في الأيام التالية.

\labelgreen{الخطوة التالية}
استخدم لوحة \textbf{الفلاتر} أعلى الشاشة لتغيير
العناصر الغذائية المعروضة (على سبيل المثال، من السعرات الحرارية إلى الكربوهيدرات) والنطاق الزمني.
انقر على علامات التبويب \textbf{أسبوعي} أو \textbf{شهري} أو \textbf{سنوي} لـ
عرض الاتجاه على مدى فترات أطول.

\suggerimento{إن مقارنة المخطط الحلقي والأهداف في
\emph{أهدافي} هي الطريقة الأكثر فعالية لتقييم ما إذا كان
نظامك الغذائي يتماشى مع أهدافك. إذا كانت شريحة واحدة أكبر بكثير من الأخريات،
يُنصح بإعادة التوازن إلى وجباتك في الأيام التالية.}

% ─── 5.3 ─────────────────────────────────────────────────────────────────────
\subsection{استخدام التقويم}

\passaggio{افتح التقويم}

\schermata[0.38\textwidth]{calendario.jpg}{شاشة \textbf{التقويم} لـ
  أبريل 2026: الأيام باللون الأحمر بدون وجبات مسجلة، واليوم 28 باللون الأخضر مع
  حد أزرق (اليوم الحالي الذي تم تسجيل وجبة فيه)، والأيام المستقبلية باللون الرمادي.
  في الأسفل توجد وسيلة الإيضاح: أخضر = وجبات مسجلة، أحمر = لا توجد وجبات.}

\labelgreen{الوصف}
انقر على أيقونة \textbf{التقويم} في شريط التنقل. رمز اللون
يحدد الأيام المسجلة على الفور: أخضر (وجبة موجودة)،
أحمر (لا توجد وجبة)، رمادي فاتح (أيام مستقبلية)، حد أزرق (اليوم الحالي).
انقر فوق يوم لعرض تفاصيل الوجبات لذلك التاريخ.

\newpage

% ════════════════════════════════════════════════════════════════════════════
%  6. إعدادات الحساب
% ════════════════════════════════════════════════════════════════════════════
\section{إعدادات الحساب}

يجمع قسم \emph{الإعدادات} جميع خيارات التخصيص:
الملف الشخصي، والإشعارات، ووحدات القياس، واللغة، وإدارة الحساب.

% ─── 6.1 ─────────────────────────────────────────────────────────────────────
\subsection{افتح الإعدادات}

\passaggio{الوصول إلى شاشة الإعدادات}

\schermata[0.38\textwidth]{impostazioni.jpg}{شاشة \textbf{الإعدادات}:
  الأقسام الثلاثة الحساب، الدعم وحول، والإجراءات مع جميع الخيارات
  المتاحة، كل منها برمز خاص به.}

\labelgreen{الوصف}
انقر على أيقونة الترس في أعلى اليمين لفتح الإعدادات.

\medskip
\begin{infobox}
  \begin{tabular}{@{}ll@{}}
    \toprule
    \textbf{الخيار}            & \textbf{الوظيفة} \\
    \midrule
    ملفي الشخصي                 & تغيير الاسم واللقب والطول وتاريخ الميلاد والجنس والصورة \\
    أهدافي                   & تحديد السعرات الحرارية المستهدفة والوزن ووحدات الماكرو والمغذيات الدقيقة \\
    الإشعارات              & تذكير بالوجبات وشرب الماء مع أوقات مخصصة \\
    اللغة والبلد       & اللغة (IT، EN، ZH، AR) والبلد المرجعي \\
    \midrule
    مساعدة                       & مركز المساعدة مع الأسئلة الشائعة \\
    سياسة الخصوصية             & وثائق حول سياسات إدارة البيانات \\
    معلومات عنا                   & معلومات حول التطبيق والفريق \\
    \midrule
    الإبلاغ عن مشكلة            & تقرير مباشر لفريق الدعم \\
    تسجيل الخروج                    & تسجيل الخروج من الحساب الحالي \\
    \bottomrule
  \end{tabular}
\end{infobox}

% ─── 6.2 ─────────────────────────────────────────────────────────────────────
\subsection{تعديل بيانات الحساب}

\passaggio{قم بتحديث ملفك الشخصي}

\schermata[0.38\textwidth]{profilo.jpg}{شاشة \textbf{ملفي الشخصي}:
  صورة الملف الشخصي مع أيقونة الكاميرا الخضراء، وحقول \emph{الاسم} (Christian)،
  \emph{اللقب} (Donzelli)، \emph{الطول} (175.0)، \emph{تاريخ الميلاد}
  (11/30/-1) وقائمة منسدلة \emph{الجنس} (ذكر). في الأسفل يوجد
  الزر الأخضر \textbf{حفظ}.}

\labelgreen{الوصف}
انقر على \textbf{ملفي الشخصي} لتعديل بياناتك الشخصية:
\begin{itemize}[leftmargin=1.4em, itemsep=2pt]
  \item \textbf{صورة الملف الشخصي} --- انقر على أيقونة الكاميرا الخضراء لالتقاط
        صورة أو اختيار واحدة من معرض الصور الخاص بك
  \item \textbf{الاسم واللقب} --- المعلومات الشخصية للمستخدم
  \item \textbf{الطول} --- بالسنتيمتر، يستخدم لحساب الاحتياجات
  \item \textbf{تاريخ الميلاد} --- يستخدم لحساب معدل الأيض الأساسي
  \item \textbf{الجنس} --- يؤثر على صيغة السعرات الحرارية
\end{itemize}

اضغط على \textbf{حفظ} للتأكيد. يقوم التطبيق تلقائيًا بتحديث الحسابات الغذائية
بناءً على البيانات الجديدة.

% ─── 6.3 ─────────────────────────────────────────────────────────────────────
\subsection{تكوين الإشعارات}

\passaggio{تمكين التذكيرات}

\dueschermate{notifiche.jpg}{شاشة \textbf{الإشعارات} مع التبديلات النشطة:
  \emph{تذكير الوجبات} مع أوقات الإفطار 7:30 صباحًا، الغداء 12:00 ظهرًا،
  العشاء 7:00 مساءً، الوجبات الخفيفة 3:00 مساءً؛ \emph{تذكير الماء} مع التكرار
  \emph{كل 30 دقيقة}.}{notifiche_orari.jpg}{قائمة \emph{تكرار التذكير}
  مفتوحة: \emph{كل 30 دقيقة} (محددة)، \emph{كل 1 ساعة}،
  \emph{كل 2 ساعة}، و \emph{أنا أقرر}. في الأسفل يوجد \emph{سجل التنبيهات}
  بأحدث الإشعارات.}

\labelgreen{الوصف}
في شاشة \emph{الإشعارات}، يتوفر مفتاحا تبديل:
\begin{itemize}[leftmargin=1.4em, itemsep=2pt]
  \item \textbf{تذكير الوجبات} --- يؤدي تشغيله إلى عرض الفترات الزمنية الأربع؛
        انقر فوق وقت أخضر لتعديله باستخدام محدد النظام
  \item \textbf{تذكير الماء} --- يؤدي تشغيله إلى عرض قائمة التكرار
        (كل 30 دقيقة، كل 1 ساعة، كل 2 ساعة، أنا أقرر)
\end{itemize}
في الأسفل يوجد \textbf{سجل التنبيهات} مع الإشعارات الأخيرة؛ انقر على
\textbf{مسح} لحذفه.

% ─── 6.4 ─────────────────────────────────────────────────────────────────────
\subsection{تغيير اللغة والبلد}

\passaggio{تكوين الموقع}

\schermata[0.38\textwidth]{lingua_paese.jpg}{شاشة \textbf{اللغة والبلد}: فتح قائمة \emph{اللغة} مع تحديد \emph{الإيطالية}، \emph{الإنجليزية}، \emph{الصينية المبسطة}، و \emph{العربية}. يوجد أسفل ذلك قسم \emph{البلد} مع \emph{إيطاليا} و \emph{إنجلترا} و \emph{الصين} و \emph{مصر}.}

\labelgreen{الوصف}
اللغات المتاحة هي الإيطالية والإنجليزية والصينية المبسطة والعربية. يتم تطبيق التغيير على الواجهة عند إعادة تشغيل التطبيق.

% ─── 6.5 ─────────────────────────────────────────────────────────────────────
\subsection{الإبلاغ عن مشكلة}

\passaggio{إرسال تقرير إلى الدعم}

\schermata[0.38\textwidth]{segnala_problema.jpg}{شاشة \textbf{الإبلاغ عن
  مشكلة}: حقل نصي بعنصر نائب \emph{صِف بالتفصيل
  المشكلة التي تمت مواجهتها\ldots} وزر أخضر \textbf{إرسال تقرير}.}

\labelgreen{الوصف}
انقر على \textbf{الإبلاغ عن مشكلة} في قسم \emph{الإجراءات}. املأ
الحقل النصي بوصف تفصيلي واضغط على \textbf{إرسال تقرير}.

\suggerimento{للحصول على حل أسرع، قم بتضمين طراز الجهاز في تقريرك
(على سبيل المثال، iPhone 14، Samsung Galaxy S23)، وإصدار
نظام التشغيل، ووصف دقيق للخطوات التي أدت إلى المشكلة.}

\bigskip
\noindent\rule{\textwidth}{1pt}{\color{nutrigreen}}
\begin{center}
  \small\itshape\color{nutrigray}--- نهاية القسم 4 ---
\end{center}

\end{document}
